\begin{exo}
	Parmi les $175$ élèves de seconde d'un lycée, il y a $56\%$ de garçons.

	Le choix de la deuxième langue vivante des élèves de seconde est résumé dans le tableau ci-dessous.

	\begin{center}
		\def\arraystretch{1.5}\small
		\begin{tabular}{|c|c|c|} \hline
			                            & \cellcolor{gray!30}Filles & \cellcolor{gray!30}Garçons \\ \hline
			\cellcolor{gray!30}Espagnol & $45,44\%$                 & $50\%$                     \\ \hline
			\cellcolor{gray!30}Allemand & $27,28\%$                 & $14,29\%$                  \\ \hline
			\cellcolor{gray!30}Italien  & $27,28\%$                 & $35,71\%$                  \\ \hline
		\end{tabular}
	\end{center}

	\begin{enumerate}
		\item Reproduire et compléter le tableau ci-dessous en indiquant les effectifs.
	\end{enumerate}
	\begin{center}
		\def\arraystretch{1.5}\small
		\begin{tabular}{|c|c|c|c|c|} \hline
			                           & \cellcolor{gray!30}Espagnol & \cellcolor{gray!30}Allemand & \cellcolor{gray!30}Italien & \cellcolor{gray!30}Total \\ \hline
			\cellcolor{gray!30}Filles  &                             &                             &                            &                          \\ \hline
			\cellcolor{gray!30}Garçons &                             &                             &                            &                          \\ \hline
			\cellcolor{gray!30}Total   &                             &                             &                            & $175$                    \\ \hline
		\end{tabular}
	\end{center}
	\begin{enumerate}[resume]
		\item On choisit un élève de seconde au hasard.
		      \begin{enumerate}
			      \item Quelle est la probabilité qu'il étudie l'espagnol ?
			      \item Quelle est la probabilité que ce soit une fille qui étudie l'allemand ?
			      \item Quelle est la probabilité que ce soit un garçon qui étudie l'italien ?
		      \end{enumerate}
	\end{enumerate}
\end{exo}
